\section{Discussion and Recommendations}
\label{section:recs}
The characteristics of the agent are described in three ways, first, the training characteristics are elaborated upon. Second, the NeuroBench results for the agents are presented. Lastly, the noise robustness is discussed. For the deployment of real closed-loop systems, these three aspects may be of varying importance. When training complex tasks, the training is a constraining factor. As discussed above, training SNN presents several additional challenges that may be detrimental. In these situations, ANN are still the algorithm of choice. Due to the challenging nature of training SNNs, sample efficiency and stability in training algorithms is a big requirement. Other reinforcement learning methods have shown to have improved efficiency and stability for training ANN. Examples of these are Proximal Policy Optimization (PPO) and Soft Actor Critic (SAC) reinforcement learning. Next to these improved algorithms, future work should try to implement SNN specific methods such as reward-modulated STDP, which has been shown to reduce the effect of dead and saturated neurons during training \cite{florian_reinforcement_2005}, in the proven RL frameworks. Furthermore, eProp has also shown interesting results \cite{eProp_pong}. 
\\
When analyzing the NeuroBench results, it can be observed that for resource constrained use cases, SNN may be a promising avenue. Their pruning opportunities allow for a significant reduction in complexity, even on hardware that is not specialized for neuromorphic algorithms. When training using RL, the SNN tends to struggle more with dead and saturated neurons, which allows for pruning without significant loss. Training ANN with similar architecture to the pruned ANN did not allow for the same performance. However, the pruning possibilities of the full ANN model were not explored. \\ 
To further analyze noise robustness of the models, more realistic noise characteristics should be injected. Currently, all inputs in the observation space receive the same white noise of mean 0 and a gain. In reality, the noise is usually smaller for sensors with a smaller range (such as the velocity). Furthermore, further work should analyze the behavior of leak and different neuron models under noise. 